% Source: https://artofproblemsolving.com/community/c3h2319138_200th_post_celebration_geometry
% Stable identity: aops-topic-2319138-q04
$\boxed{\text{False}}$. A ray is named by its endpoint first, followed by any other point of the ray. The endpoint of this ray is $D$, so its names are $\overrightarrow{DP}$ and $\overrightarrow{DT}$.

\begin{asy}
usepackage("amsmath");
size(5.5cm);
defaultpen(fontsize(10pt));
pair D=(0,0), P=(2.5,0), T=(5,0);
draw(D--(6.5,0),linewidth(0.9),Arrow(6));
dot(D); dot(P); dot(T);
label("$D$",D,S); label("$P$",P,S); label("$T$",T,S);
label("endpoint",D,N,fontsize(8pt));
\end{asy}

The name $\overrightarrow{TP}$ would denote the ray starting at $T$ and going through $P$, which points the other way (towards $D$ and beyond), and $\overrightarrow{PD}$ would denote the ray starting at $P$ and going through $D$; neither is the given ray.
